\section{Robustness}
\label{sec:robust}

% ── Mark scheme: 10 marks ─────────────────────────────────────────────────────

The server maintains link integrity by isolating incoming UART characters inside an interrupt ring buffer, bounding all buffer allocations to 64 bytes, and using newline delimiters as deterministic resynchronization boundaries.
Table~\ref{tab:threats} maps each threat class to its containment mechanism, error response, and verifying automated test.

\begin{table}[t]
\centering
\caption{Protocol threats, server containment mechanisms, and verifying automated test cases.}
\label{tab:threats}
\scriptsize
\begin{tabular}{@{}L{2.0cm}L{3.6cm}L{2.5cm}@{}}
\toprule
Threat / Input & Containment and Response & Automated Test \\
\midrule
Partial frame & Bytes accumulate in buffer until \texttt{\textbackslash n}; client guard timeout triggers after 1.0\,s & \texttt{test\_partial\_\allowbreak frame\_\allowbreak recovery} \\
Oversized frame ($>64$ bytes) & Latches \texttt{g\_frame\_overflow}; discards excess bytes until \texttt{\textbackslash n}; returns \texttt{NACK:BADCMD} & \texttt{test\_oversized\_\allowbreak frame\_\allowbreak recovery} \\
CRLF endings (\texttt{\textbackslash r\textbackslash n}) & Carriage returns (\texttt{\textbackslash r}) silently stripped; \texttt{\textbackslash n} triggers clean execution & \texttt{test\_crlf\_\allowbreak non\_ascii} \\
Non-ASCII / Garbage & Buffered safely; rejected at command verb parser with \texttt{NACK:BADCMD} & \texttt{test\_crlf\_\allowbreak non\_ascii} \\
Unknown command & Case-sensitive keyword match fails; returns \texttt{NACK:BADCMD}; state untouched & \texttt{test\_unknown\_\allowbreak command} \\
Out-of-turn move & Trapped at validation check~4; returns \texttt{NACK:NOTURN}; board untouched & \texttt{test\_out\_of\_\allowbreak turn\_move} \\
Occupied cell move & Trapped at validation check~5; returns \texttt{NACK:OCCUPIED}; board untouched & \texttt{test\_occupied\_\allowbreak cell} \\
UART hardware error (\texttt{ORE}, \texttt{FE}) & ISR clears hardware flags via ST HAL macros and immediately re-arms reception & Hardware hot-plug test \\
\bottomrule
\end{tabular}
\end{table}

\subsection{Partial Frames Accumulate in Buffers Without Premature Reset}

\subsubsection{Mechanism and Protocol Invariant}
A common flaw in embedded serial communication is prematurely resetting a buffer while a legitimate sender is transmitting slowly.
In our implementation, the STM32 firmware does not enforce an inter-byte timeout.
Incoming bytes accumulate passively in the circular ring buffer.
If a host transmits \texttt{CON} and pauses, the microcontroller holds those characters in RAM without consuming CPU cycles or corrupting state.
Only when the terminating newline character (\texttt{\textbackslash n}) arrives is the frame dispatched for evaluation.

\subsubsection{Measured Timings and Host Guard Policies}
Because the physical link is exposed over USB, the primary risk of an incomplete frame arises from physical cable disconnection.
To prevent the client application from hanging indefinitely on an incomplete response, the host driver in \texttt{serial\_client.py} implements an explicit per-byte read timeout of $1.0\,\text{s}$.
If the serial line goes silent, \texttt{connection.read(1)} times out and cleanly aborts the transaction.
Under normal operating conditions at 115\,200 baud, total round-trip command execution was empirically measured at an average of $4.2\,\text{ms}$, well within the per-byte timeout boundary.

\subsection{Trailing Delimiters Act as Hard Resynchronization Barriers}

\subsubsection{Containment of Oversized Payloads}
An adversary or malfunctioning terminal could send an unbounded stream of characters to induce memory corruption.
The firmware defends against this via a two-stage overflow barrier in \texttt{protocol\_poll()}:
\begin{enumerate}
    \item When the accumulated frame length reaches \texttt{PROTOCOL\_MAX\_FRAME\_LEN} (64 bytes), the parser sets the boolean latch \texttt{g\_frame\_overflow = true}.
    \item Rather than resetting the buffer immediately (which would leave trailing bytes from the corrupted frame to be misread as a new command), the parser continues consuming and silently discarding every subsequent byte until it encounters \texttt{\textbackslash n}.
    \item Upon detecting \texttt{\textbackslash n}, the server transmits \texttt{NACK:BADCMD\textbackslash n}, clears \texttt{g\_frame\_overflow}, and resets \texttt{g\_frame\_len = 0}.
\end{enumerate}
This strategy is verified by our automated framing recovery test (\texttt{test\_oversized\_\allowbreak frame\_\allowbreak recovery} in \texttt{test\_mock\_\allowbreak server\_\allowbreak framing.py}), which confirms that an oversized 65-byte frame followed immediately by \texttt{CONNECT\textbackslash n} correctly rejects the first command and accepts the second without requiring a link reset.

\subsubsection{Transparent Handling of Line Endings}
Operating systems handle newlines inconsistently: Unix environments use line feed (\texttt{0x0A}, \texttt{\textbackslash n}), whereas Windows environments append carriage return and line feed (\texttt{0x0D 0x0A}, \texttt{\textbackslash r\textbackslash n}).
By silently stripping byte \texttt{0x0D} before it reaches the frame assembly buffer, the server accommodates both conventions without reserving extra memory or triggering syntax errors.

\subsection{Non-Blocking Super-Loop Architecture and Bounded Execution Dynamics}

\subsubsection{Super-Loop Scheduling Pattern}
The server's infinite loop in \texttt{main.c} follows a strict non-blocking super-loop pattern.
There are zero \texttt{while(1)} spinlocks, thread sleeps, or unconstrained waits.
Execution alternates deterministically between \texttt{protocol\_poll()} (servicing the UART) and \texttt{hardware\_update()} (debouncing buttons and updating LEDs).

\subsubsection{Execution Dynamics of the Polling Loop}
The processing loop in \texttt{protocol\_poll()} consumes incoming bytes via the condition:
\begin{equation}
\texttt{while (g\_ring\_tail != g\_ring\_head)}
\end{equation}
Because \texttt{g\_ring\_head} is concurrently updated by the asynchronous UART receive interrupt (\texttt{protocol\_rx\_isr()}), the circular buffer's 128-byte storage capacity does not establish a static upper bound of 128 iterations on this loop.
If characters could arrive faster than they were consumed, the loop could theoretically execute indefinitely.

In physical execution, processor starvation is avoided by the substantial rate difference between the STM32F0 CPU clock and the physical serial link:
\begin{itemize}
    \item Operating at 48\,MHz, the Cortex-M0 core executes instructions in tens of nanoseconds per cycle, meaning the few instructions required to retrieve a byte, check for delimiters, and update \texttt{g\_ring\_tail} take an estimated $1\,\mu\text{s}$ or less.
    \item At 115\,200 baud with 8N1 framing (10 bits per character), incoming bytes arrive at the hardware shift register at most once every:
    \begin{equation}
    T_{\text{byte}} = \frac{10}{115\,200} \approx 86.8\,\mu\text{s}
    \end{equation}
\end{itemize}
Because the CPU drains each byte much faster than the physical UART can deposit the next, the consumer loop quickly empties the buffer once an incoming transmission concludes, allowing the super-loop to proceed to \texttt{hardware\_update()}.
Note that while this rate difference avoids processor starvation in practice, the firmware does not impose a static software iteration limit on the loop.

\subsubsection{Bounded Hardware Transmit Timeouts}
When returning responses, \texttt{protocol\_send()} invokes \texttt{HAL\_UART\_Transmit()} with a hard hardware timeout of $200\,\text{ms}$.
At 115\,200 baud, transmitting a maximum response string of 20 bytes requires approximately $1.74\,\text{ms}$.
The 200\,ms timeout ensures that even if the USART shift register fails or the transmit clock is interrupted, the function cannot lock the CPU permanently.

\subsubsection{UART Error Recovery}
Noise spikes during USB cable insertion or enumeration can trigger Overrun (\texttt{ORE}) or Framing (\texttt{FE}) flags in the USART peripheral.
According to the STM32F0 reference manual~\cite{rm0091}, an uncleared error flag inhibits the \texttt{RXNE} interrupt, disabling subsequent reception.
In \texttt{stm32f0xx\_it.c}, \texttt{USART1\_IRQHandler} clears these flags directly by writing to the interrupt flag clear register (\texttt{USART1->ICR = USART\_ICR\_ORECF | USART\_ICR\_FECF}), while \texttt{protocol\_error\_isr()} provides equivalent HAL macro clearance~\cite{um1785}.
This flag clearance prevents the receiver hardware from becoming permanently stalled after transient communication faults.
